% =========================================================
%  Capítulo 3 — Metodologia
% =========================================================

\chapter{Metodologia}
\label{cap:metodologia}

Este capítulo descreve a metodologia adotada para o planejamento, desenvolvimento
e validação da plataforma.

% -----------------------------------------------------------
\section{Classificação da Pesquisa}
\label{sec:classificacao-pesquisa}

Do ponto de vista da sua natureza, este trabalho caracteriza-se como uma pesquisa
aplicada, pois tem como objetivo gerar conhecimento para solução de um problema
específico \cite{[REFERENCIA_METODOLOGIA]}. Quanto à abordagem, é predominantemente
qualitativa, uma vez que a avaliação do sistema envolve análise de funcionalidades
e usabilidade. Com relação aos procedimentos técnicos, trata-se de um estudo de
caso de desenvolvimento de \textit{software}.

% -----------------------------------------------------------
\section{Metodologia de Desenvolvimento}
\label{sec:metodologia-desenvolvimento}

O desenvolvimento foi conduzido de forma iterativa e incremental, com ênfase
em entregas funcionais a cada ciclo. As funcionalidades foram organizadas em
etapas:

\begin{enumerate}
  \item \textbf{Modelagem do domínio:} levantamento e especificação das entidades,
        regras de negócio e casos de uso;

  \item \textbf{Implementação do \textit{backend}:} desenvolvimento em Java/Spring Boot
        seguindo a Arquitetura Hexagonal, com cobertura de testes unitários;

  \item \textbf{Integração e testes:} validação dos componentes por meio de testes
        unitários da camada de domínio (modelos e serviços) e testes de controladores;

  \item \textbf{Documentação:} geração automática da documentação da API via Swagger/
        OpenAPI e elaboração deste documento.
\end{enumerate}

% -----------------------------------------------------------
\section{Ferramentas e Ambiente de Desenvolvimento}
\label{sec:ferramentas}

O Quadro~\ref{qdr:ferramentas} lista as principais ferramentas e tecnologias
utilizadas no desenvolvimento da plataforma.

\begin{table}[H]
  \centering
  \caption{Ferramentas e tecnologias utilizadas}
  \label{qdr:ferramentas}
  \begin{tabular}{lll}
    \toprule
    \textbf{Categoria}        & \textbf{Tecnologia/Ferramenta}    & \textbf{Versão} \\
    \midrule
    Linguagem                 & Java                              & 17              \\
    Framework                 & Spring Boot                       & 3.5.6           \\
    Persistência              & Spring Data JPA / Hibernate       & 6.x             \\
    Banco de dados            & PostgreSQL                        & 15              \\
    Mapeamento objeto-relacional & MapStruct                      & 1.5.x           \\
    Autenticação              & JWT (JJWT)                        & 0.12.x          \\
    Armazenamento             & AWS S3 (SDK v2)                   & 2.x             \\
    Conteinerização           & Docker / Docker Compose           & 24.x            \\
    Build                     & Apache Maven                      & 3.x (wrapper)   \\
    Documentação de API       & Springdoc OpenAPI (Swagger UI)    & 2.x             \\
    IDE                       & IntelliJ IDEA                     & --              \\
    Controle de versão        & Git / GitHub                      & --              \\
    Testes unitários          & JUnit 5 + Mockito                 & --              \\
    \bottomrule
  \end{tabular}
  \fonte{Elaborado pelo autor (\imprimirdata).}
\end{table}

% -----------------------------------------------------------
\section{Estratégia de Testes}
\label{sec:estrategia-testes}

A estratégia de testes adotada concentra-se na camada de domínio, em conformidade
com os princípios da Arquitetura Hexagonal, que recomenda que o núcleo da aplicação
seja testável de forma isolada, sem dependência de frameworks ou infraestrutura.

Os testes foram implementados com JUnit 5 e Mockito, cobrindo:

\begin{itemize}
  \item \textbf{Testes de modelos de domínio:} verificação das regras de negócio
        encapsuladas nos métodos das entidades, como \texttt{User.awardXp()} e
        os construtores via \textit{factory methods};

  \item \textbf{Testes de serviços (casos de uso):} verificação do comportamento
        dos \texttt{ServiceImpl} com repositórios \textit{mockados}, garantindo
        o fluxo correto de orquestração entre domínio e portas de saída;

  \item \textbf{Testes de controladores:} verificação dos \textit{endpoints} REST
        com \texttt{@WebMvcTest}, cobrindo casos de sucesso, validações de entrada
        e tratamento de erros HTTP.
\end{itemize}

Testes de integração (\texttt{@SpringBootTest}, \texttt{@DataJpaTest}) estão fora
do escopo deste trabalho.

% -----------------------------------------------------------
\section{Controle de Versão}
\label{sec:controle-versao}

O código-fonte foi versionado com Git, hospedado em repositório privado no GitHub.
O fluxo de trabalho adotado foi baseado em \textit{commits} incrementais organizados
por funcionalidade, com mensagens no formato convencional
(\textit{feat}, \textit{fix}, \textit{refactor}).
